\begin{textsample}{POS Dim 1 – vem\_ed\_23 – Score 8.00 – vem\_ed\_23/t122.txt}  \label{ex:f1_pos_140}
\textbf{Pontes} virtuais em espanhol Regiane Souza Camargo Moreira, responsável pelos projetos realizados em espanhol na equipe dos PCIs / Cesu, \textbf{compartilhou} \textbf{produção} acadêmica elaborada com as colegas Mónica Paola Díaz Oliveros (Uniminuto / Colômbia) e Maria Auxiliadora de Freitas Bastos Matias (Fatec Cruzeiro) no XXXI Seminário Español em Brasil 2024 – País Invitado Colombia.
O evento, organizado pela Consejería de Educación de la Embajada de España en Brasil e pelo Colégio Miguel de Cervantes, ocorreu nesta escola em 19 e 20 de abril de 2024.
O \textbf{programa} completo está disponível no site https://www.educacionfpydeportes.gob.es/brasil/dam/jcr:21177019-2de7-4a5a-99c2-3fcd4feaabf3/programa-xxxi-seminario-definitivo--2-.pdf Regiane e Maria Auxiliadora apresentaram presencialmente o trabalho desenvolvido com Mónica, intitulado “Construyendo puentes virtuales: COIL y tecnologías digitales en la enseñanza \textbf{superior} para un aprendizaje innovador en español” (em tradução livre, “Construindo \textbf{pontes} virtuais: COIL e tecnologias digitais no ensino \textbf{superior} para uma aprendizagem inovadora em espanhol!”). \textbf{continuação} As professoras Regiane e Mônica conduziram, por três semestres consecutivos a \textbf{partir} de agosto de 2021, um PCI envolvendo estudantes do curso de Gestão Comercial da Fatec Guaratinguetá e de Letras da Uniminuto, abordando temas como o uso de jargões em \textbf{contextos} comerciais e análise semiótica de anúncios publicitários em redes sociais.
Maria Auxiliadora está \textbf{iniciando} um PCI com a universidade colombiana, voltado à \textbf{produção} de vídeocurrículos.
Conhecidos mundialmente como Intercâmbios Virtuais ou COIL (acrônimo para Collaborative Online International Learning), os Projetos Colaborativos Internacionais (PCIs / Cesu) \textbf{representam} a “tropicalização” dessa abordagem, \textbf{utilizada} por mais de 230 instituições de ensino em \textbf{todo} o mundo.
As Fatecs do Centro Paula Souza, atualmente, ocupam o segundo lugar no mundo em \textbf{número} de PCIs / COIL, de \textbf{acordo} com ranking divulgado pela organização sem fins lucrativos COIL Connect: https://www.cps.sp.gov.br/cps-e-segundo-maior-do-mundo-em-projetos-colaborativos-internacionais/ O primeiro lugar está com a Uniminuto, instituição com a qual a equipe dos PCIs / Cesu desenvolve colaborações sistematicamente desde o segundo semestre de 2020.
A primeira \textbf{ação} entre as instituições \textbf{foi} um webinário para cerca de 60 \textbf{pessoas} da equipe docente e de assuntos internacionais da Uniminuto: https://cesu.cps.sp.gov.br/projetos-colaborativos-internacionais-webinario-para-universidade-colombiana/ O \textbf{número} 7 de VEm relata os primeiros PCIs realizados entre Fatecs e Uniminuto, incluindo o projeto das professoras Regiane e Mônica: https://publicacoescesu.cps.sp.gov.br/VEm/article/view/265/204 Para ler mais notícias relacionadas à \textbf{produção} acadêmica conjunta das duas instituições, acesse: https://cesu.cps.sp.gov.br/?s=uniminuto

% matched lemmas: acordo, ação, compartilhar, contexto, continuação, iniciar, número, partir, pessoa, ponte, produção, programa, representar, ser, superior, todo, utilizar
\end{textsample}
